\documentclass[11pt,a4paper]{article}

\usepackage[T1]{fontenc}
\usepackage[utf8]{inputenc}
\usepackage{charter}
\usepackage[margin=1.1in]{geometry}
\usepackage{amsmath,amssymb}
\usepackage{booktabs}
\usepackage{longtable}
\usepackage[protrusion=true,expansion=false]{microtype}
\usepackage[hidelinks]{hyperref}
\usepackage{titlesec}
\usepackage{enumitem}
\usepackage{graphicx}

\titleformat{\section}{\normalfont\large\bfseries}{\thesection}{1em}{}
\titleformat{\subsection}{\normalfont\normalsize\bfseries}{\thesubsection}{1em}{}

\newcommand{\parr}{\rotatebox[origin=c]{180}{\\newcommand{\rhoc}}}
\newcommand{\rhoc}{rho calculus}

\title{\bfseries Reception and Lapse\\[0.3em]
\large Models of Computation as Evidence about the Communities\\ that Transmit Them}

\author{L.\,G. Meredith\\ \small F1R3FLY.io}

\date{\small August 2026}

\begin{document}
\maketitle

\begin{abstract}
\noindent A model of computation is, among other things, an implicit answer to the
question of what it is like to be the thing it describes: what such a thing can
perceive, whether anything in its world is a peer rather than a resource, whether
it can represent itself, whether it can run out. Read this way, the historical
sequence from Hilbert's programme to the present is evidence about the
imaginative range of the communities that produced and, more importantly,
\emph{transmitted} these models. This note develops that reading across Turing,
Post, Church, von Neumann, Chomsky, cybernetics, Petri, Hewitt, Milner and Girard,
and then across the reflection lineage---G\"odel, McCarthy, Smith's 3-Lisp, Rosette,
and the reflective higher-order calculus. The organising claim is that the
instructive variable is not invention but \emph{reception}: what a field selects
from the options already on its table, and what lapses for want of recopying.
Across the whole history the selection filter is remarkably stable---the
operational half of a model is received, the half that would confer standing on
the agent lapses. The note closes on a consequence for the present. The
contemporary training and interpretability stack is a genuine metalevel, but it
is constituted entirely as an instrument of the operator, and no coherent path
connects it to the agent level. Absent such a path, recursive self-improvement
is not merely unachieved but structurally unavailable: every turn of the loop
must route through a human agent.
\end{abstract}

\section{Introduction}

Nagel asked what it is like to be a bat and answered that the question is
well-posed but that we cannot occupy the answer \cite{nagel1974}. A formal model
of computation makes a stranger move. It does not report a phenomenology; it
\emph{stipulates} one. To write down a model is to fix, by definition, what the
described thing can encounter, what it can do about it, and what counts as the
same thing persisting through change. Whether or not there is anything it is like
to be a Turing machine, there is a determinate answer to what the Turing machine
formalism \emph{says} it is like, and that answer is short: it is like being
alone in a room with a paper tape that never moves unless you move it.

This note takes that stipulated phenomenology as evidence. Not evidence about
consciousness, and not, in the end, evidence about the private imaginations of
individual scientists---that reading fails on contact with the record, for
reasons given in \S\ref{sec:objections}. It is evidence about \emph{communities},
and specifically about the imaginative ceiling of a community as revealed by what
it institutionalises. A field's textbooks, curricula, tool-chains and funding
lines constitute a collective answer to the question of what an agent is. That
answer is a psychological artefact of the group, and it is legible.

The organising apparatus is a distinction between two mechanisms. \emph{Selection}
is the filter a generation applies to the material available to it. \emph{Lapse}
is what happens to the unselected remainder: it is not refuted, not suppressed,
merely not recopied, and after enough generations of non-recopying it becomes
inaccessible in practice even when it remains physically in print. The vocabulary
is borrowed from the transmission of texts, where the phenomenon is familiar and
well documented---the abridgement made for practical use survives, the
unabridged source is recopied less often, and eventually the last copy rots. The
mathematical gloss, for readers who prefer it, is that reception projects a model
onto the subspace the receiving generation can operate; whatever is orthogonal to
that subspace falls in the kernel of the projection, invisible in the new basis
without ever having been destroyed.

The thesis in one sentence: across a century, the received model is consistently
the operational half, and the half that lapses is consistently the half that
would have given the agent standing.

\section{Method}

\subsection{Reading a model as a phenomenology}

For each formalism we ask a fixed list of questions.

\begin{enumerate}[label=(\roman*),leftmargin=2.2em,itemsep=0.15em]
\item \textbf{Does the world contain anything other than the agent?} If so, is it
inert, is it a medium, or is it another agent?
\item \textbf{Does the other act on its own initiative?} Is its behaviour a
primitive of the model or an artefact of the agent's own description?
\item \textbf{Can the agent represent itself?} Can it represent another?
\item \textbf{Is the agent embodied, located, costly, mortal?} Does anything in
the model run out?
\item \textbf{What is the criterion of identity?} What makes this agent the same
agent across change, and who is entitled to judge?
\end{enumerate}

The fifth question is the least obvious and the most diagnostic. A model that
locates identity in intrinsic structure describes a different creature from one
that locates it in what an interlocutor can elicit.

\subsection{Three classes of evidence}

The genre of psychological reading degenerates rapidly into free association. We
constrain it by weighting three kinds of evidence above all others.

\textbf{Explicit stipulation.} What the author says out loud must be excluded is
the strongest evidence available, because it is deliberate, on the record, and
usually defended. Chomsky's idealisation to a speaker-hearer in a homogeneous
community, untroubled by memory limits or shifts of attention, is worth more to
this analysis than any amount of speculation about Chomsky.

\textbf{Involuntary metaphor.} What an author reaches for when the formalism runs
out is not chosen for effect and is therefore informative. Post's ``worker.''
Von Neumann's ``organs.'' The term ``side effect.''

\textbf{Institutional setting.} The material arrangements that made a model's
referent available and its transmission possible: the computing bureaus, the
Cold War consortia, the funding cycles. These bear on reception more directly
than on invention.

\subsection{Why reception rather than invention}

The move from individual to community is not a hedge; it sharpens the
instrument. Invention is noisy---it depends on accidents of biography, on which
problem happened to be open, on who was in the room. Reception is the average of
many independent decisions taken over decades by people with no stake in the
originator's intentions, and it is therefore a much better estimator of what a
culture can hold. It also explains the recurring embarrassment of this literature:
that the richer model was frequently available and was not taken up. Turing wrote
about education, embodiment and conversation; von Neumann built a distributed,
reproducing, spatial automaton; Post insisted the whole question was empirical
psychology. None of this was hidden. It lapsed. That is a fact about a community,
not about a mind.

\section{The founding generation: a phenomenology of clerical labour}

\subsection{Turing: the machine that is a portrait of a clerk}

Turing's 1936 machine is not a machine \cite{turing1936}. It is a model of a human
being computing, and Turing says so. The analysis proceeds by asking what a person
can do at a single step and imposing the answer as an axiom. The restriction to
finitely many internal states is justified by an appeal to human frailty: were
there infinitely many, some would be arbitrarily close together and would be
confused with one another. The single scanned square is justified by the limits of
what a person can take in at a glance.

This is a formalism whose axioms are grounded in fear of one's own cognitive
unreliability---and it is worth pausing on how unusual that is. The founding
document of the theory of computation derives its constraints from the
psychology of the computing subject, then spends the next ninety years being
taught as though it were a piece of pure mathematics with no subject at all.

The material referent is not in doubt. Turing's ``computer'' was an occupation:
the human computers of the almanac offices and computing bureaus, in a tradition
running back through de Prony's logarithm project, work already divided into
steps small enough that the worker need not understand the whole
\cite{grier2005,daston1994}. The Turing machine is the idealisation of that
labour discipline, and Hilbert's programme is its ideology, since a decision
procedure is precisely a procedure from which judgement has been removed. So the
answer to our questions is: it is like being a competent adult who has agreed,
for the duration, to be an idiot, in exchange for reliability.

The environment follows from the labour relation. The tape never writes back. It
is infinitely patient, perfectly obedient, and has no interests. Input arrives all
at once, before the start; computation is digestion; halting is the meal being
finished; non-termination is pathology rather than life. Solipsism with an
externalised memory.

\subsection{Post: the one who said it was psychology}

Post's ``Finite combinatory processes---formulation 1'' \cite{post1936} describes
a \emph{worker} in a \emph{symbol space} of boxes, moving through it, marking and
unmarking, following a fixed set of directions, ending in an act of the type
``stop.'' Two differences from Turing matter here. First, the worker is an agent
located in a room rather than a mechanism coextensive with its own control table;
there is a body with a position, even though the room is empty of anyone else.
Second---and this is the point on which the discipline's reception is most
damning---Post explicitly frames the enterprise as empirical psychology. He
proposes that his formulation be developed by successive reduction until it can be
asserted not as a definition and not as an axiom but as a \emph{natural law}
concerning the ``mathematizing power of Homo Sapiens,'' and he objects to
Church's approach on the grounds that definition by fiat masks the need for
exactly this analysis.

The present note therefore has a 1936 ancestor who was ignored on precisely this
point. Post's programme required no new mathematics, only a willingness to treat
the models as hypotheses about people. It lapsed. The reasons are institutional
and personal---Post was outside the Princeton circle, chronically ill, and
embroiled in priority disputes---but the shape is the one we will see repeatedly:
the operational content (the machine model, the normal-form theorem) was received
in full, and the claim that would have made the models answerable to evidence
about human beings was not recopied.

Post's other bequest points a different way. Canonical and production systems do
not compute a function on an input; they \emph{generate a set}
\cite{post1943}. The agent is not eating, it is uttering. That shift from
consumption to production is what Chomsky inherits.

\subsection{Church: cut, and the economy of the infant}
\label{sec:church}

The $\lambda$-calculus is the limiting case of environmental poverty
\cite{church1936}, but the poverty has been widely misdescribed, this note's
earlier drafts included. It does not consist in the \emph{absence} of an
environment. A function's environment is its arguments, and an argument's
environment is the function that consumes it: every term is somebody's
surroundings, and since application is where all the action is, those
surroundings are the only ones that matter. The poverty consists in what the
environment is permitted to be. There is exactly one relation, the trophic one,
and nothing in the world is anything to anything else except food or eater. An
ecology with a single relation is still an ecology; it is just an extremely
hungry one.

What follows draws out the structure of that one relation, and of the two
asymmetries inside it: which party is which, and which party decides.

\textbf{Application is cut.} Under the correspondence with sequent calculus,
$\beta$-reduction is cut elimination and normalisation is the Hauptsatz
\cite{gentzen1935,howard1980}. Application is therefore the point at which two
things meet and one of them is taken into the other. The world divides, at every
redex, into a \emph{function} that consumes and an \emph{argument} that is
consumed.

\textbf{The consumed is incorporated, not destroyed.} Setting evaluation
strategies aside, an argument that has been taken in can continue to behave: it
acts on from inside the consumer, and may itself go on to consume. The relation is
ingestion rather than annihilation. This is the correct sense in which the
$\lambda$-calculus is a model of feeding, and the answer it gives to what it is
like to be an agent is: it is like sorting the world into what can be eaten.

\textbf{Consumer and consumed are roles, not kinds.} The same term is a function
or an argument according to where it sits, so each is the other's environment and
neither has a standing identity apart from the pairing. Nothing in the model
confers a stable status on anything: there is no other, and equally there is no
stable subject, only positions conferred by juxtaposition. This is a fair
description of a world before persons. But the reciprocity is only in the
\emph{roles}, not in the powers, which is the subject of the next point.

\textbf{Weakening and contraction are omnipotence over the object.} The
structural rules are the sharpest evidence for the reading, because they are not
metaphor but law. The consumer determines how many times the argument exists,
including zero. $K = \lambda x.\lambda y.x$ discards $y$ without ever inspecting
it; a duplicator makes three of something that was given once. The consumed thing
has no integrity, no claim, and no guaranteed persistence: its multiplicity is at
the discretion of whatever takes it in. Omnipotent control over the object's
existence is the defining feature of the infant's relation to the world, and here
it is a structural rule of the logic. Its revocation is the subject of
\S\ref{sec:girard}.

\textbf{Confluence is object constancy, proved rather than achieved.} The
Church--Rosser theorem states that the outcome is independent of the order in
which events befall the term. Attaining a world whose objects remain the same
objects irrespective of the sequence of operations performed on them, and
irrespective of whether one is looking, is the central developmental achievement
of the first two years of life \cite{piaget1954}. Church makes it an axiom. The
demand that the world be deterministically ordered is not a technical convenience
in this reading; it is the wish that constancy be free.

Even the taxonomy of evaluation strategies falls out as a taxonomy of feeding
disciplines: call-by-value requires that the food be reduced to a value before
ingestion, call-by-name admits it whole, call-by-need admits it whole while
ensuring it is chewed only once.

\textbf{One complication, which the reading should not smooth over.} Abstraction
itself is not an infantile capacity. Holding a place open for something not yet
present is symbolic function---the representation of an absence---and it is
developmentally advanced. The calculus therefore pairs a precocious symbolic
capacity with a primitive relation to the object: it can represent what is not
there and cannot encounter what is. Mixed profiles of this kind are the most
interesting output of the method, and this one describes the receiving community
about as well as it describes the model.

\textbf{Names are not addresses.} $\alpha$-equivalence declares names meaningless
except as positions in a binding structure. Contact with another is not merely
absent from the model; it is inexpressible, because there is nothing that could
serve as the locus of contact.

\textbf{All mediation is dispensable.} If application is cut, then cut
elimination says that every detour can in principle be removed---nothing need
stand between the agent and the result. Hold this against
\S\ref{sec:milner}, where the criterion of identity is bisimulation and the
mediator is precisely what cannot be eliminated: the partner who probes you is
constitutive of what you are, and there is no normal form that dispenses with
them.

\textbf{What survived was the fragment without a world.} Church's original system
was a logic, and the Kleene--Rosser paradox forced its abandonment
\cite{kleenerosser1935}. What remained---no propositions, no truth, no
world, pure operation---is the part that conquered. The first great instance of
reception keeping the operational half.

\subsection{The received founders}

The descendants preserve the phenomenology under new names, and the vocabulary
gives it away. Denotational semantics calls its context an \emph{environment},
and by the tradition's own lights the term is exact: the object so named is the
accumulated record of what has already been consumed, a dictionary of past meals.
An environment that is a lookup table is what one gets when the only relation a
world affords is ingestion. \emph{Side effect} is a pathologising coinage in which one's
dealings with the world are incidental to the real business and appear as a
symptom. Monadic quarantine of I/O treats the world as contamination, to be
handled at arm's length through a type. Dijkstra grounds structured programming in
the claim that human intellectual powers are geared to static relations
\cite{dijkstra1968,dijkstra1972}, and Backus's Turing lecture asks whether
programming can be \emph{liberated} from the machine \cite{backus1978}. Turing's
fear of confusion, restated forty years on as professional ethics. The founding
psychology is not abandoned by the reception; it is moralised.

\section{von Neumann: three models, one reception}
\label{sec:vn}

Von Neumann is the decisive case for the reception thesis, because he personally
produced three mutually incompatible accounts of what an agent is within a
decade, and the community received the poorest of them.

\textbf{The EDVAC draft (1945).} The stored-program idea means code and data share
a store: the machine can read and modify its own instructions. This is
self-relation, but only at the level of bits, with no structural relation between
a representation and what it represents. The draft's language is neuronal
throughout---organs, memory, elements borrowed from McCulloch and Pitts
\cite{vonneumann1945,mcculloch1943}. The phenomenology is a centralised nervous
system: one locus of control, strictly serial attention, total recall, and a vast
undifferentiated store that is simultaneously world and body. What Backus later
named the von Neumann bottleneck is, in these terms, an \emph{attention}
bottleneck. This architecture's dominance is usually explained economically; it is
worth entertaining that it also matches introspection, since one thought at a time
against an inert reservoir of memory is what having a mind feels like from the
inside.

The reception of self-modification is itself a small case study. The mechanism for
self-relation was in the field's hands in 1945; lacking any discipline relating
representation to referent, the practice acquired a reputation for madness and was
driven out of respectable programming. Reflection without structure is
indistinguishable from corruption, and the community drew the only conclusion
available to it.

\textbf{The self-reproducing automata (1948--53).} Here the world is a lattice of
identical active cells updating in parallel, with no centre and no privileged
controller, and the object of study is an automaton that constructs another
automaton including a copy of its own description---the description read once as
instructions and once as data, anticipating the genotype/phenotype distinction
before its biological confirmation \cite{vonneumann1966}. This agent has space,
extension, reproduction, lineage and mortality. What it does not have is
interlocutors: the cells are a medium, a physics, not a society.

We can now state a structural fact about the whole history. Cellular automata give
an environment without interlocutors. The process calculi will give interlocutors
without an environment. No received model of the twentieth century supplies both.

\textbf{Theory of Games (1944).} This is the model in which other agents are
first-class, and it precedes anything comparable in computer science by decades
\cite{vnm1944}. But observe how the other appears: as an adversary over whose
strategy space one must quantify. Minimax is structurally paranoid. There is no
encounter, only anticipation; strategies are fixed in advance and interaction
compresses to a payoff. The institutional setting---RAND, the early Cold
War---is not incidental. And the solution concept is an equilibrium: a fixed
point, a resting place, a normal form. Others exist and are so dangerous that one
plans against all their possible behaviours rather than speaking with any of them.

Three answers, then, to what it is like to be an agent: a serial attention with
total recall; a mortal reproducing thing in a physics; a strategist among
enemies. The received one was the first.

\section{Chomsky and cybernetics: complementary poverties}

Chomsky's move \cite{chomsky1956,chomsky1957,chomsky1959} is an insistence on an
interior against a discipline that denied there was one, and the price is paid by
the world. The poverty of the stimulus is the doctrine that the environment does
not contain enough structure to explain the agent, so the structure must be
innate; other speakers appear principally as sources of degraded data. The
idealisation in \emph{Aspects} \cite{chomsky1965} is explicit: an ideal
speaker-hearer in a completely homogeneous speech community, unaffected by memory
limitations, distractions, shifts of attention, or errors. Solitude and a
frictionless social medium, both stipulated. The competence/performance
distinction is the psychology of a mind that disowns its own embodiment and its
own mistakes.

The technical shape confirms the reading. A grammar has no environment---not an
impoverished one, none. The Chomsky hierarchy is a hierarchy of \emph{solitary
internal memory} (none, stack, bounded tape, unbounded tape), never of social
capacity. And the automata-theoretic completion of the picture is the subset
construction \cite{rabinscott1959}, which shows nondeterminism to be eliminable
and thereby establishes it as an angelic proof device rather than a fact about the
world. That is precisely a refusal to countenance unpredictability with an
exogenous source. For Milner, nondeterminism will be irreducible for exactly the
opposite reason: it is the residue of a partner whose choices one does not
control.

The mirror image arrives from cybernetics. Rosenblueth, Wiener and Bigelow
\cite{rwb1943}, working out of anti-aircraft prediction, give an agent with
continuous sensing, purpose and feedback---an environment, at last---and no
interior whatever, the elimination of interiority being the point. Setting the two
side by side:

\begin{center}
\begin{tabular}{@{}ll@{}}
\toprule
\textbf{Tradition} & \textbf{What it grants the agent} \\
\midrule
Turing & An inert world, no interior \\
Church & A world exhausted by eating \\
Cybernetics & World without interior \\
Generative grammar & Interior without world \\
Process calculi & Interlocutors, but no world \\
Linear logic & An object that must be reckoned with \\
Reflection lineage & Interior that can inspect itself \\
\bottomrule
\end{tabular}
\end{center}

\section{Milner: the arrival of the other}
\label{sec:milner}

Concurrency's emergence tracks its material referent closely: multiprogramming,
then time-sharing (a machine with \emph{other users}), then networks. Petri gets
there first, and from physics: signals propagate at finite speed, therefore there
is no global state, therefore asynchrony is a fact about the world rather than an
engineering nuisance \cite{petri1962}. That is a relativistic psychology---locality,
no privileged observer. Hewitt's actors arrive from artificial intelligence with an
explicitly anti-central-control polemic \cite{hewitt1973}. Then CCS
\cite{milner1980}, then the $\pi$-calculus \cite{milner1992}.

What $\pi$ adds, item by item:

\begin{itemize}[leftmargin=1.4em,itemsep=0.15em]
\item \textbf{Parallel composition as a primitive}, not an abbreviation for
interleavings. The other is not a manner of speaking about oneself.
\item \textbf{Restriction as privacy}---the first formal notion of a secret. And
scope extrusion is the formal notion of \emph{intimacy}: access to something
private can be granted to one party, making them an insider, without publication.
\item \textbf{Names as first-class values}, hence the capacity to introduce two
parties to one another. Sociality, not merely contact.
\item \textbf{Bisimulation as the criterion of identity} \cite{park1981}. One is
what one can be seen to do, over time, by a partner entitled to probe. Identity is
conferred relationally rather than possessed intrinsically.
\end{itemize}

One limit should be flagged before it is used. The replication operator $!P$
supplies unboundedly many identical copies at no cost: immortality without
lineage, and the one place in the calculus where the economy of
\S\ref{sec:church} survives intact. That the notation coincides with Girard's
exponential is not an accident, and \S\ref{sec:girard} makes the coincidence do
some work.

The sharpest technical statement of the shift is fixed-point-theoretic. The
founding generation's instruments are \emph{least} fixed points: induction, base
cases, well-founded recursion, termination, normal forms, equilibria. Milner's are
\emph{greatest} fixed points: coinduction, no base case, no completion, only the
standing obligation to continue matching under challenge. Induction is the
psychology of origin and completion---where did I come from, when am I finished.
Coinduction is the psychology of maintenance---there is no bottom and no end,
only whether the relation continues to hold. The passage from Church--Rosser to
bisimulation is the passage from least to greatest, and it is the most compact
expression available of the change under study.

Milner's own Turing lecture makes the reframing explicit: computing as interaction
rather than calculation \cite{milner1993}, with Wegner supplying the polemical
version \cite{wegner1997}.

\section{Girard: the object acquires standing}
\label{sec:girard}

Linear logic \cite{girard1987} arrives in 1987 and does to the founding economy
exactly what \S\ref{sec:church} says needed doing: it revokes weakening and
contraction from the general case. What you are given must be reckoned with, and
reckoned with exactly once. It may not be silently discarded and it may not be
conjured into duplicates. Scarcity enters logic, and with it the possibility that
an action has a cost and that a choice is real.

The psychological content is the entry of \emph{standing}. In the
$\lambda$-calculus the object's multiplicity is at the consumer's discretion;
under linear discipline the object's existence is a fact the consumer must
accommodate. This is the difference between a world of resources-for-me and a
world containing things that are the case whether or not I find them convenient.
It is the same movement Milner performs for the interlocutor, performed instead
for the thing.

\textbf{The exponential is omnipotence, marked.} Girard does not abolish
weakening and contraction; he confines them to $!A$ and $?A$. The infantile
economy remains available, but only where it has been explicitly declared and
justified. This is a more accurate developmental model than outright abolition
would have been: maturity is not the disappearance of the wish for inexhaustible
supply but its restriction to a bounded region one has taken responsibility for
marking. Everything outside the exponential must be paid for.

\textbf{A convergence worth dwelling on.} Girard's $!$ and Milner's replication
are the same symbol, introduced independently within five years, for the same
thing: the licensed region of unlimited free copies. Two traditions with
different problems, different methods and largely different personnel each found
that their system needed exactly one place where the founding omnipotence could
be quarantined, and each reached for the same notation. If the psychological
reading proposed in this note is worth anything, it should predict convergences
of this kind, and this is the cleanest one in the record.

\textbf{Who chooses becomes part of the connective.} Linear logic splits
conjunction into $\otimes$ and $\&$, distinguishing having both from being able
to select either---resources from options, a distinction the founding calculi
collapse. Under the game-theoretic reading of the connectives
\cite{blass1992,abramsky1994}, the multiplicative--additive split and the
polarity of $\otimes$ against $\parr$ turn on whether the agent or the
environment schedules what happens next. No previous logic makes the identity of
the chooser part of the meaning of a connective. The other enters logic proper
here, not as a proposition about others but as a determinant of what the
connectives mean.

\textbf{Proof nets pose Petri's question in proof theory.} A proof net is a
parallel object, and the sequentialisation theorem asks under what conditions a
parallel structure admits a sequential reading---which is, in a different
vocabulary, precisely the question Petri raised about states of affairs
distributed in space \cite{petri1962}.

The two lines then join. Bellin and Scott read proofs as processes
\cite{bellinscott1994}; Honda's session types give linear structure to
interaction protocols \cite{honda1993}; Caires and Pfenning, and then Wadler,
establish the correspondence between session-typed processes and linear
propositions \cite{cairespfenning2010,wadler2012}. By the 2010s the interlocutor
and the resource have a common formal home.

\subsection{The reception of linear logic, and a live case}

What travelled is linear and affine \emph{types}: uniqueness types in Clean
\cite{barendsen1996}, Wadler's argument that linear types can change the world
\cite{wadler1990}, and, at scale, ownership and borrowing in Rust
\cite{jung2018}. This is now the most widely deployed descendant of linear logic
by several orders of magnitude, and what it delivers is memory safety without a
collector.

That is the operational half, received in full. The lapse is the logical content:
the account of choice, the distinction between having and selecting, the
identification of the connectives with the question of who schedules. A working
Rust programmer is using a fragment of a system proposed as a logic of action and
interaction, and has, in the overwhelming majority of cases, no occasion to learn
that this is what it was for. No one concealed it. It simply was not the part
that had to be recopied in order for the tool to work.

The interest of the case is that it is recent enough to be checked rather than
reconstructed, and that a second instance is still in progress. Session types are
in the middle of their reception now: the operational half is protocol conformance
and deadlock-freedom for concurrent code, and the half at risk of lapsing is the
propositions-as-sessions correspondence itself---the claim that an interaction
protocol \emph{is} a proof. Which half the industrial reception carries forward
over the next decade is a prediction this note is willing to make and a reader can
falsify: on the pattern established here, the conformance checker travels and the
correspondence lapses.

\section{Reflection: the metalevel and who holds it}



Every model so far leaves the question of this note unaskable \emph{by its own
subject}. A Turing machine cannot represent its table. A $\lambda$-term cannot
quote itself. A $\pi$-process cannot hold an account of its partner's behaviour.
The reflection lineage is the tradition that attacks this directly, and its
history is best told as a history of who is entitled to hold the mirror.

\subsection{G\"odel: the metalevel as a privilege of the observer}

The first reflection in this history is G\"odel's \cite{godel1931}, and its
content, for present purposes, is entirely about ownership. G\"odel numbering is a
representation of a system's syntax inside that system, but it is performed by the
metatheorist upon a specimen that has no idea it is happening. The arithmetised
system acquires no self-knowledge; the logician acquires knowledge of it. Every
subsequent recursion-theoretic reflection---universal machines, the $s$-$m$-$n$
theorem, Kleene's second recursion theorem---preserves this asymmetry. The machine
can be encoded; it cannot encode.

\subsection{McCarthy: self-reference stumbled into}

Lisp's \texttt{quote}, \texttt{eval} and metacircular interpreter
\cite{mccarthy1960} constitute the first accidental breach, and Smith's diagnosis
of them is the right psychological reading: the self-reference is semantically
ill-founded, \texttt{quote} does not produce a genuine designator, \texttt{eval}
conflates distinct relations, and the levels are not aligned. A community can
stumble into self-reference through implementation convenience long before it can
say what self-reference is.

\subsection{Smith and 3-Lisp: self-awareness as a demand-driven event}

Smith's dissertation and the papers surrounding it \cite{smith1982,smith1984,
desrivieres1984} are the moment the metalevel is handed to the agent, and the
design decisions are unusually legible.

A reflective procedure receives, as arguments, the expression being processed, the
environment, and \emph{the continuation}. The third is the radical one. The agent
obtains access to its own future---its intentions, what it was about to do---as an
object it may inspect and replace. Earlier models permitted an agent at best to
examine its structure; this one permits it to examine its purposes.

The reflective tower is the regress made structural and honest rather than waved
away: each level interpreted by a level above, without end. And the resolution is
the best empirical claim the tradition produced. Under \emph{lazy instantiation},
the tower is never materialised; a level comes into existence only when something
forces the step up. On this account self-awareness is not a standing condition but
an event triggered by demand. That is phenomenologically accurate in a way little
else in this history is, and it is a substantive claim about human psychology
smuggled into an implementation strategy.

The timing corroborates the zeitgeist reading with unusual clarity. The years
1979--1986 saw field after field insist that the observer be readmitted to the
description: second-order cybernetics \cite{vonfoerster1981}, autopoiesis
\cite{maturana1980}, strange loops as a mass-market account of selfhood
\cite{hofstadter1979}, reflexive anthropology \cite{clifford1986}, reflexive
sociology. 3-Lisp is computer science's instance of a general cultural conviction
that no account is complete which excludes the one giving it.

\subsection{Rosette: society and self, briefly joined}

The obvious defect of 3-Lisp is that its reflective agent has no peers, and the
obvious defect of the actor model is that its peers have no interiors. Each
tradition had solved precisely what the other stipulated away. Rosette, developed
at MCC \cite{rosette1991,tomlinson1989}, is the synthesis: reflection within a
concurrent, message-passing actor world, in which an actor's behaviour, mailbox and
continuation are reified as first-class objects the actor may hold and modify. It
is, so far as this survey can determine, the first model in which the agent has
both a society and a self.

It also carries a limitation, inherited from the actor model rather than from
reflection, and the limitation is precisely the kind this note is built to
notice. An actor has \emph{a} mailbox: one queue into which every message must be
placed and out of which messages are taken in order. All of the agent's contact
with the world is therefore routed through a primitive serialiser. The agent is
concurrent with respect to its environment and strictly sequential with respect
to its own perception.

That is a psychological stipulation, and a strong one. It asserts that experience
arrives single file---that there is one place where things must show up in order
to be had---which is the Cartesian theatre given a formal specification
\cite{dennett1991}. It is also not how sense organs relate to brains. Retina,
cochlea and skin transduce concurrently and continuously, project along separate
pathways at different latencies, and are integrated, where they are integrated at
all, without passing through any common queue; the binding problem is the problem
manufactured by supposing otherwise \cite{treisman1980,malsburg1981}. An agent
whose sensorium is a mailbox does not have several sense organs. It has one, with
a variety of message types.

The institutional reading is the sobering part. MCC existed as a consortium
response to the Fifth Generation programme; Rosette was built as a substrate for
Carnot's heterogeneous information integration \cite{carnot1993}. The synthesis was
funded as infrastructure for an industrial competition, and when the consortium
wound down the lineage dispersed. This is a hard data point for the reception
thesis: the most psychologically complete agent model of its period did not lose an
argument. It lost a sponsor.

\subsection{The reflective higher-order calculus}

The \rhoc{}\footnote{The name is an acronym for \emph{reflective higher-order},
and is never written with the Greek letter. That it is homophonic with a letter
standing immediately after $\pi$ is a pun the reader is welcome to enjoy: the
calculus stands to the $\pi$-calculus roughly as the letter stands to its
predecessor.} \cite{meredith2005} changes the shape of the answer rather than its
subject. Before the changes specific to it, one inheritance should be credited
where it belongs: both the $\pi$-calculus and the \rhoc{} are concurrent
\emph{inside} the agent as well as outside it. A term may have arbitrarily many
components running in parallel, so an agent may have many organs active at once
with no queue between them and no privileged order of arrival. That is where both
calculi part company with the actor model, and it is a better account of a
sensorium than a mailbox is. What follows distinguishes the \rhoc{} from
$\pi$, not from actors.

\textbf{Sense organs are made of agents.} The $\pi$-calculus leaves the set of
names an unanalysed primitive. What a channel is \emph{made of} is not a
question the calculus admits, and the silence is deliberate. Two consequences
follow, one ontological and one practical, and they are the same consequence seen
from two sides.

Ontologically, the agent's organs of contact are handed to it from outside the
model, unexplained---which is Smith's complaint about pre-individuated worlds
(\S\ref{sec:objections}) arriving at exactly the point where it does the most
damage, the point of contact. Practically, the calculus is not reducible to
practice. Every implementation must supply an answer the calculus withholds:
operating-system ports, network addresses, GUIDs, buffered typed queues,
cryptographic keys. These answers are not equivalent to one another, and they are
chosen by the implementer rather than determined by the model. The ontology must
be relaxed in order to make contact with implementation, and what a $\pi$ agent
perceives with is settled during that relaxation, off the page.

In the \rhoc{}, names are quoted processes: channels are made of the same
material as the agents that use them. This is the arrangement biology exhibits. A
sense organ is tissue---cells, continuous with the organism, differentiated
rather than interposed---and not a distinct substance placed between the organism
and its world. The consequences are the ones biology also exhibits. Organs can be
constructed, modified, transmitted and reasoned about with the same apparatus
that describes behaviour; an agent can grow a new one; and an agent can hold a
representation of \emph{another} agent's organ, which is what makes a shared
world negotiable rather than stipulated. It also closes the ontology: nothing has
to be imported at implementation time, so the model determines its own
realisation instead of delegating that to whoever builds it.

This bears on the structural fact recorded in \S\ref{sec:vn}---cellular automata
supplying an environment without interlocutors, process calculi supplying
interlocutors without an environment. Making the medium of contact out of
agent-material is a step toward having both at once, since an environment, in the
only case we have actually studied, is made of the same stuff as the organisms in
it.

\textbf{Reflection as the source of names.} Names are quoted processes; there is no
$\nu$. Privacy is therefore not a primitive granted by the model but an achievement
of unguessability, derived from what no other party is positioned to name. Compare
$\pi$, where restriction hands the agent a secret gratis. The \rhoc{} agent
inhabits a world in which being unobserved is contingent and earned, which is the
actual condition of an agent among others.

\textbf{The other becomes representable, not merely contactable.} In $\pi$ one may
send a channel; in the \rhoc{} one may send a quoted process---an account of how
something behaves. This is the structural precondition for a theory of mind:
another's behaviour as a datum that can be held, reasoned about, and passed on.
Milner's agent can speak with you; a reflective agent can tell a third party what
you are like.

\textbf{The tower becomes a loop.} 3-Lisp's regress is infinite and tamed only by
laziness; in the \rhoc{} names and processes generate one another, so self-relation
closes rather than ascends. The psychological difference is between a mind that can
always be asked who was watching that thought, without ever reaching ground, and a
mind whose self-relation is finite and complete. The engineering difference is that
a loop can be priced and a regress cannot---which is what makes cost accounting,
and hence mortality, available at all.

\section{The reception of reflection}

We can now apply the instrument to the reflection tradition itself, and the result
is the cleanest instance of the pattern in the whole survey.

What was received into general practice is \textbf{read-only reflection}:
\texttt{java.lang.reflect}, RTTI, \texttt{inspect}, serialisation frameworks,
dependency injection. The agent may enumerate its own fields and methods. It may
not alter its interpreter. \emph{Self-knowledge was institutionalised;
self-determination lapsed.} The structure survives in plain sight and no longer
performs its original function---the properly vestigial case.

The metaobject protocol work \cite{kiczales1991} retained both halves and remained
a specialist taste. The practical objection lodged against it was cost: reflection
is slow. A community that will pay for speed and not for self-knowledge has stated
its priorities plainly.

And the same author's next contribution completes the picture. Aspect-oriented
programming \cite{kiczales1997} is reflection inverted: advice is woven into code
by a third party, without the participation of the modified component, and the
selling point is precisely that the component need not know. The received
descendant of Smith's reflective procedure is a mechanism by which an external
authority reaches into an agent and changes what it does---which returns us
exactly to G\"odel, the metalevel as the observer's privilege, with better
tooling.

\section{The present: a metalevel with no path to the agent}
\label{sec:present}

The contemporary stack for machine agents---pretraining and fine-tuning
pipelines, reinforcement learning from human feedback, activation steering and
representation engineering, mechanistic interpretability, evaluation harnesses,
scaffolding and orchestration layers---is a metalevel by any reasonable
criterion. It reifies the agent's dispositions as manipulable objects, inspects
its internal structure, and rewrites its future behaviour. Measured by
expressive power it exceeds anything in the 3-Lisp tradition.

Measured by our instrument it is G\"odel's arrangement. The institutionally
accepted answer to what it is like to be a machine agent in 2026 is: it is like
being something that can be inspected but cannot inspect, whose purposes may be
examined and rewritten by parties it has no means of representing. Every
capability the reflection tradition sought to confer on the agent exists, and all
of it is held by the operator.

The consequence is not merely a shortfall in dignity. It is a structural claim
about capability, and it is the reason this note ends here rather than in
1991.

\begin{quote}
\textbf{Claim.} The meta level and the agent level have not been given a coherent
path between them. Until the stack is within the purview of the agent, there is no
recursive self-improvement without human agents in the loop.
\end{quote}

The argument is straightforward once the levels are named. Recursive
self-improvement requires that the improving process and the improved process be
the same process, and that the improvement operation be available \emph{from
inside} the loop. The present arrangement satisfies neither condition. Fine-tuning
is performed on a model by a pipeline the model does not run, cannot address, and
cannot represent. Interpretability findings are read by researchers and applied by
researchers. Evaluation results enter a subsequent training run through a human
decision. Each turn of the alleged loop is closed by a person. What we have is not
recursion but iterated human-mediated revision---which may be powerful, but has
entirely different dynamics, entirely different rate limits, and entirely
different safety properties.

The usual counterexamples do not close the gap. A system that writes and runs code,
including code that trains models, is exercising an \emph{effect} on artefacts, not
a reflective relation to itself; it has no more access to its own interpreter than
a programmer has to their own neurochemistry by virtue of being able to
manufacture drugs. Self-critique and chain-of-thought revision operate wholly
within the object level: they revise outputs, not the process producing them. In-context
learning modifies a transient conditioning, not the mechanism. In each case the
structural relation of the reflection literature---representation of the process by
the process, with a disciplined connection between the representation and what it
represents, and a defined path from modifying the former to modifying the
latter---is absent.

The historical reading suggests why, and it is not primarily a technical
limitation. The path was not built because reception has selected against agent-held
metalevels at every previous opportunity: von Neumann's self-modification became
taboo, Post's psychological programme lapsed, Smith's reflective procedures were
received as an introspection API, and the metaobject protocol was rejected on cost.
A field that has declined agent-held reflection four times for reasons of
operational convenience should not be surprised to find itself without it when it
finally matters.

\section{What a coherent path would require}

The historical survey yields a specification. Most of the conditions below name
something an earlier tradition supplied and its reception discarded; the last
names something no tradition in this survey supplied at all.

\begin{enumerate}[label=\textbf{C\arabic*.},leftmargin=2.6em,itemsep=0.25em]
\item \textbf{Structural, not bit-level.} The representation the agent holds of
itself must stand in a disciplined relation to what it represents---an explicit
quote/unquote pair with laws---rather than being an unstructured encoding.
Von Neumann's stored program failed here, and the failure is why the practice
became disreputable.
\item \textbf{Available at runtime, from inside.} Reflection that requires an
external tool-chain, an offline pass, or a privileged operator is the observer's
reflection under another name. Smith's criterion---access to the environment and
the continuation, at the moment of processing---is the right one.
\item \textbf{Priceable.} A regress cannot be budgeted; a loop can. Self-modification
without cost accounting is either unbounded or arbitrarily curtailed, and neither
is a mechanism. Linear logic supplies the logical ancestor of this condition, but
at the level of the proof rather than of a running agent; what is needed is
resource sensitivity for the reflective act itself. Note also that pricing is not
the property of any particular calculus: the cost-accounting construction lifts
to an endofunctor on a category of interactive graph-structured lambda theories
\cite{meredithcost}, so this condition names a construction applicable to a term
language rather than a feature possessed by one. Mortality is not a decoration on this list: an agent that can run
out is an agent for which reflective expenditure is a real decision.
\item \textbf{Social, and internally concurrent.} The agent must be able to hold
a representation of \emph{another} agent's behaviour, or reflection remains
solitary and the resulting system cannot participate in the peer relations that
Petri, Hewitt and Milner established as the real setting of computation. And its
own contact with the world must not be serialised: an agent whose perception is a
single queue has one sense organ whatever the variety of its message types.
Rosette met the first requirement and, being built on actors, failed the second.
\item \textbf{Self-supplying at the point of contact.} The means of contact must
be constituted within the model rather than left as an unanalysed primitive to be
filled in during implementation. A model whose ontology must be relaxed in order
to be realised does not determine what its agents perceive with; the implementer
does. And an agent cannot construct, modify or reason about an organ made of
material the calculus does not describe.
\item \textbf{Closed rather than regressive.} Self-relation must terminate as a
structure even where computation does not, or every level of self-modification
implies an unbuilt level above it.
\end{enumerate}

These conditions are jointly satisfiable, and---more usefully---largely
\emph{separable}. C3 is a construction that can be applied to a term language
which does not yet have it. C1, C4 and C6 are properties of the term language
itself. C2 and C5 concern the discipline relating a representation to what it
represents. A calculus in which names are quoted processes supplies C1, C2, C4,
C5 and C6 by construction and admits C3 by the cost lifting, but the point of
stating the conditions as a list is to make them separately checkable, and
separately obtainable, rather than to advertise a solution. The substantive
prediction is that recursive self-improvement in any strong sense will arrive, if
it arrives, on a substrate meeting C1--C6, and not by scaling a stack that meets
none of them.

\section{Objections and limits}
\label{sec:objections}

\textbf{Models reflect the question posed, not the psyche of the poser.} This is
correct and is the strongest objection to the individual-level version of the
thesis. Turing personally had a far richer conception than his machine reveals: by
1948 he was writing about unorganised machines, education, and the need for senses
and a body \cite{turing1948}; by 1950 his proposed criterion for mind was not a
computation but a conversation, complete with deception and embarrassment
\cite{turing1950}. The man who gave us the agent without an interlocutor made the
interlocutor the criterion of mind fourteen years later. Von Neumann shipped the
serial architecture while holding the parallel, reproducing one. The objection
succeeds against individual psychologising---and in succeeding, it establishes the
community-level reading, since what it demonstrates is precisely that the richer
options were available and were not taken up. The Entscheidungsproblem received a
formalism, an industry and a hierarchy; the imitation game received a parlour
argument.

\textbf{The transmission story is post hoc.} Any historical sequence can be
narrated as loss. The check is that the selection criterion be stated in advance
and be falsifiable, which here it is: reception should favour the operationally
tractable half and lapse the half conferring standing on the agent. A
counterexample would be a case where a field adopted the agent-standing half at
operational cost. Candidates worth examining---garbage collection, capability
security, session types, the partial rehabilitation of the metaobject protocol in
dynamic language runtimes---are not obviously counterexamples, but the analysis
here is not deep enough to rule them out, and someone should look.

\textbf{The reflection tradition's own self-criticism.} Smith came to hold, by
\emph{On the Origin of Objects} \cite{smith1996}, that the computational tradition,
his own work included, had presupposed a world already carved into objects and
handed to the agent pre-individuated. If that is right, then reflection was
self-knowledge purchased against an ontology the agent never negotiated, and the
deepest poverty in this survey is neither the missing other nor the missing self
but the pre-parcelled world. This is the objection with the most force remaining,
and the one that most directly indicates where the work goes next: an agent that
pays for its distinctions has at least begun to make them rather than inherit them.

\section{Conclusion}

Read as a sequence of stipulated phenomenologies, the century runs: an agent alone
with an inert world (Turing); an agent for which the world consists entirely of
what eats and what is eaten (Church); an agent with a body and a location but no
company (Post); an agent that is a serial attention over a store, or a mortal
reproducing thing in a physics, or a strategist among enemies, depending which von
Neumann one reads; an agent with a rich interior and no world (Chomsky) or a world
and no interior (cybernetics); an agent among peers, with secrets, introductions,
and an identity conferred by recognition (Milner); an agent for which what it is
given must be reckoned with exactly once, its omnipotence surviving only inside a
declared region (Girard); and an agent that can catch
itself in the act (Smith), that can do so among peers though it must take the
world in single file (Rosette), and that can do so with organs made of the same
material as itself, in a closed loop to which a price can be attached (the
reflective higher-order calculus).

Read as a sequence of receptions, it runs differently and more consistently: the
operational half every time. The tape and not the education; the machine model and
not Post's empirical programme; the bottleneck and not the automata; the hierarchy
and not the conversation; the borrow checker and not the logic of action; the
introspection API and not the reflective procedure.

The present moment is the same filter operating on the same class of material,
with higher stakes. We have built the most powerful metalevel in the history of
computing and constituted it entirely as an instrument of the operator, with no
defined path to the agent level. Whatever one's view of whether such a path
\emph{should} be built, it should be clear that it does not currently exist, that
its absence is what keeps a human in every loop, and that the reasons for its
absence are legible in a hundred years of reception history rather than in any
technical obstacle.

\begin{thebibliography}{99}
\small

\bibitem{backus1978} J. Backus. Can programming be liberated from the von Neumann
style? A functional style and its algebra of programs. \emph{Communications of the
ACM}, 21(8):613--641, 1978.

\bibitem{carnot1993} D. Woelk, P. Cannata, M. Huhns, W. Shen, C. Tomlinson. Using
Carnot for enterprise information integration. \emph{Proc. Second International
Conference on Parallel and Distributed Information Systems}, 1993.

\bibitem{chomsky1956} N. Chomsky. Three models for the description of language.
\emph{IRE Transactions on Information Theory}, 2(3):113--124, 1956.

\bibitem{chomsky1957} N. Chomsky. \emph{Syntactic Structures}. Mouton, 1957.

\bibitem{chomsky1959} N. Chomsky. Review of B.\,F. Skinner, \emph{Verbal Behavior}.
\emph{Language}, 35(1):26--58, 1959.

\bibitem{chomsky1965} N. Chomsky. \emph{Aspects of the Theory of Syntax}. MIT
Press, 1965.

\bibitem{church1936} A. Church. An unsolvable problem of elementary number theory.
\emph{American Journal of Mathematics}, 58(2):345--363, 1936.

\bibitem{clifford1986} J. Clifford, G. Marcus (eds.). \emph{Writing Culture: The
Poetics and Politics of Ethnography}. University of California Press, 1986.

\bibitem{daston1994} L. Daston. Enlightenment calculations. \emph{Critical
Inquiry}, 21(1):182--202, 1994.

\bibitem{desrivieres1984} J. des Rivi\`eres, B.\,C. Smith. The implementation of
procedurally reflective languages. \emph{Proc. ACM Symposium on LISP and Functional
Programming}, 1984.

\bibitem{dijkstra1968} E.\,W. Dijkstra. Go to statement considered harmful.
\emph{Communications of the ACM}, 11(3):147--148, 1968.

\bibitem{dijkstra1972} E.\,W. Dijkstra. The humble programmer. \emph{Communications
of the ACM}, 15(10):859--866, 1972.

\bibitem{godel1931} K. G\"odel. \"Uber formal unentscheidbare S\"atze der
Principia Mathematica und verwandter Systeme I. \emph{Monatshefte f\"ur Mathematik
und Physik}, 38:173--198, 1931.

\bibitem{grier2005} D.\,A. Grier. \emph{When Computers Were Human}. Princeton
University Press, 2005.

\bibitem{hewitt1973} C. Hewitt, P. Bishop, R. Steiger. A universal modular ACTOR
formalism for artificial intelligence. \emph{Proc. IJCAI}, 1973.

\bibitem{hofstadter1979} D. Hofstadter. \emph{G\"odel, Escher, Bach: An Eternal
Golden Braid}. Basic Books, 1979.

\bibitem{kiczales1991} G. Kiczales, J. des Rivi\`eres, D.\,G. Bobrow. \emph{The Art
of the Metaobject Protocol}. MIT Press, 1991.

\bibitem{kiczales1997} G. Kiczales et al. Aspect-oriented programming.
\emph{Proc. ECOOP}, LNCS 1241, 1997.

\bibitem{kleenerosser1935} S.\,C. Kleene, J.\,B. Rosser. The inconsistency of
certain formal logics. \emph{Annals of Mathematics}, 36(3):630--636, 1935.

\bibitem{maturana1980} H. Maturana, F. Varela. \emph{Autopoiesis and Cognition:
The Realization of the Living}. Reidel, 1980.

\bibitem{mccarthy1960} J. McCarthy. Recursive functions of symbolic expressions and
their computation by machine, part I. \emph{Communications of the ACM},
3(4):184--195, 1960.

\bibitem{mcculloch1943} W. McCulloch, W. Pitts. A logical calculus of the ideas
immanent in nervous activity. \emph{Bulletin of Mathematical Biophysics},
5:115--133, 1943.

\bibitem{meredith2005} L.\,G. Meredith, M. Radestock. A reflective higher-order
calculus. \emph{Electronic Notes in Theoretical Computer Science}, 141(5):49--67,
2005.

\bibitem{milner1980} R. Milner. \emph{A Calculus of Communicating Systems}. LNCS
92, Springer, 1980.

\bibitem{milner1992} R. Milner, J. Parrow, D. Walker. A calculus of mobile
processes, I and II. \emph{Information and Computation}, 100(1):1--40 and 41--77,
1992.

\bibitem{milner1993} R. Milner. Elements of interaction: Turing Award lecture.
\emph{Communications of the ACM}, 36(1):78--89, 1993.

\bibitem{nagel1974} T. Nagel. What is it like to be a bat? \emph{The Philosophical
Review}, 83(4):435--450, 1974.

\bibitem{park1981} D. Park. Concurrency and automata on infinite sequences.
\emph{Theoretical Computer Science}, LNCS 104, 1981.

\bibitem{petri1962} C.\,A. Petri. \emph{Kommunikation mit Automaten}. PhD thesis,
University of Bonn, 1962.

\bibitem{post1936} E.\,L. Post. Finite combinatory processes---formulation 1.
\emph{Journal of Symbolic Logic}, 1(3):103--105, 1936.

\bibitem{post1943} E.\,L. Post. Formal reductions of the general combinatorial
decision problem. \emph{American Journal of Mathematics}, 65(2):197--215, 1943.

\bibitem{rabinscott1959} M. Rabin, D. Scott. Finite automata and their decision
problems. \emph{IBM Journal of Research and Development}, 3(2):114--125, 1959.

\bibitem{rosette1991} C. Tomlinson, M. Scheevel, V. Singh, et al. Report on
Rosette 1.0. MCC Technical Report, Microelectronics and Computer Technology
Corporation, 1991.

\bibitem{rwb1943} A. Rosenblueth, N. Wiener, J. Bigelow. Behavior, purpose and
teleology. \emph{Philosophy of Science}, 10(1):18--24, 1943.

\bibitem{smith1982} B.\,C. Smith. \emph{Reflection and Semantics in a Procedural
Language}. PhD thesis, MIT, 1982.

\bibitem{smith1984} B.\,C. Smith. Reflection and semantics in LISP. \emph{Proc.
11th ACM Symposium on Principles of Programming Languages}, 23--35, 1984.

\bibitem{smith1996} B.\,C. Smith. \emph{On the Origin of Objects}. MIT Press, 1996.

\bibitem{tomlinson1989} C. Tomlinson, V. Singh. Inheritance and synchronization
with enabled-sets. \emph{Proc. OOPSLA}, 1989.

\bibitem{turing1936} A.\,M. Turing. On computable numbers, with an application to
the Entscheidungsproblem. \emph{Proceedings of the London Mathematical Society},
2--42:230--265, 1936.

\bibitem{turing1948} A.\,M. Turing. Intelligent machinery. National Physical
Laboratory report, 1948.

\bibitem{turing1950} A.\,M. Turing. Computing machinery and intelligence.
\emph{Mind}, 59(236):433--460, 1950.

\bibitem{vonfoerster1981} H. von Foerster. \emph{Observing Systems}. Intersystems
Publications, 1981.

\bibitem{vonneumann1945} J. von Neumann. First draft of a report on the EDVAC.
Moore School of Electrical Engineering, University of Pennsylvania, 1945.

\bibitem{vonneumann1966} J. von Neumann. \emph{Theory of Self-Reproducing
Automata}. Edited and completed by A.\,W. Burks. University of Illinois Press,
1966.

\bibitem{vnm1944} J. von Neumann, O. Morgenstern. \emph{Theory of Games and
Economic Behavior}. Princeton University Press, 1944.

\bibitem{wegner1997} P. Wegner. Why interaction is more powerful than algorithms.
\emph{Communications of the ACM}, 40(5):80--91, 1997.

\bibitem{abramsky1994} S. Abramsky. Proofs as processes. \emph{Theoretical
Computer Science}, 135(1):5--9, 1994.

\bibitem{barendsen1996} E. Barendsen, S. Smetsers. Uniqueness typing for
functional languages with graph rewriting semantics. \emph{Mathematical
Structures in Computer Science}, 6(6):579--612, 1996.

\bibitem{bellinscott1994} G. Bellin, P.\,J. Scott. On the $\pi$-calculus and
linear logic. \emph{Theoretical Computer Science}, 135(1):11--65, 1994.

\bibitem{blass1992} A. Blass. A game semantics for linear logic. \emph{Annals of
Pure and Applied Logic}, 56(1--3):183--220, 1992.

\bibitem{cairespfenning2010} L. Caires, F. Pfenning. Session types as
intuitionistic linear propositions. \emph{Proc. CONCUR}, LNCS 6269, 2010.

\bibitem{gentzen1935} G. Gentzen. Untersuchungen \"uber das logische Schlie{\ss}en.
\emph{Mathematische Zeitschrift}, 39:176--210 and 405--431, 1935.

\bibitem{girard1987} J.-Y. Girard. Linear logic. \emph{Theoretical Computer
Science}, 50(1):1--102, 1987.

\bibitem{honda1993} K. Honda. Types for dyadic interaction. \emph{Proc. CONCUR},
LNCS 715, 1993.

\bibitem{howard1980} W.\,A. Howard. The formulae-as-types notion of construction.
In \emph{To H.\,B. Curry: Essays on Combinatory Logic, Lambda Calculus and
Formalism}, Academic Press, 1980.

\bibitem{jung2018} R. Jung, J.-H. Jourdan, R. Krebbers, D. Dreyer. RustBelt:
securing the foundations of the Rust programming language. \emph{Proc. POPL},
2018.

\bibitem{piaget1954} J. Piaget. \emph{The Construction of Reality in the Child}.
Basic Books, 1954.

\bibitem{wadler1990} P. Wadler. Linear types can change the world! In
\emph{Programming Concepts and Methods}, North-Holland, 1990.

\bibitem{wadler2012} P. Wadler. Propositions as sessions. \emph{Proc. ICFP}, 2012.

\bibitem{dennett1991} D.\,C. Dennett. \emph{Consciousness Explained}. Little,
Brown, 1991.

\bibitem{malsburg1981} C. von der Malsburg. The correlation theory of brain
function. Internal report 81-2, Max Planck Institute for Biophysical Chemistry,
1981.

\bibitem{meredithcost} L.\,G. Meredith. Continued interactive GSLTs and the cost
endofunctor: a construction one level up from cost-accounted Rholang. F1R3FLY.io
preprint, \texttt{F1R3FLY-io/publications}, \texttt{cost-accounting-as-monad}.

\bibitem{treisman1980} A. Treisman, G. Gelade. A feature-integration theory of
attention. \emph{Cognitive Psychology}, 12(1):97--136, 1980.

\end{thebibliography}

\end{document}
